\documentclass[12pt]{article}
\usepackage{amsmath}
\usepackage{graphicx}
\usepackage{float}
\usepackage[utf8]{inputenc}
\usepackage[T1]{fontenc}
\usepackage[english]{babel}
\usepackage[shortlabels]{enumitem}
\usepackage{empheq}
\usepackage{multicol}
\usepackage{microtype}
\usepackage{msc}
\usepackage[skip=1ex]{subcaption}
\usepackage{nomencl}
\usepackage{fancyhdr}
\usepackage{hyperref}
\usepackage{url}
\usepackage{xcolor}
\usepackage{listings} % Added for beautiful JSON formatting
\usepackage{tabularx}
\usepackage{xurl}
\usepackage{amssymb}
\usepackage[strings]{underscore}
\usepackage{fvextra}
\usepackage{csquotes}

\hypersetup{
      colorlinks=true,
      linkcolor=blue!60!black,
      citecolor=green!50!black,
      urlcolor=blue!70!black
}


\pagestyle{fancy}
\fancyheadoffset{2cm}

\newcommand{\monthnameonly}{%
\ifcase\month\or January\or February\or March\or April\or May\or June\or July\or August\or September\or October\or November\or December\fi
}



\setlength {\marginparwidth }{2cm}
\newcommand\JSONnumberstyle{\color{blue}}
\newcommand\JSONstringstyle{\color{red}}
\newcommand\JSONlabelstyle{\color{black}\bfseries}


\lstdefinelanguage{json}{
  basicstyle=\ttfamily\small,
  showstringspaces=false,
  breaklines=true,
  backgroundcolor=\color{gray!10},
  literate=
   *{0}{{{\JSONnumberstyle0}}}{1}
    {1}{{{\JSONnumberstyle1}}}{1}
    {2}{{{\JSONnumberstyle2}}}{1}
    {3}{{{\JSONnumberstyle3}}}{1}
    {4}{{{\JSONnumberstyle4}}}{1}
    {5}{{{\JSONnumberstyle5}}}{1}
    {6}{{{\JSONnumberstyle6}}}{1}
    {7}{{{\JSONnumberstyle7}}}{1}
    {8}{{{\JSONnumberstyle8}}}{1}
    {9}{{{\JSONnumberstyle9}}}{1}
    {:}{{{\color{black}{:}}}}{1}
    {,}{{{\color{black}{,}}}}{1}
    {\{}{{{\color{black}{\{}}}}{1}
    {\}}{{{\color{black}{\}}}}}{1}
    {[}{{{\color{black}{[}}}}{1}
    {]}{{{\color{black}{]}}}}{1},
  stringstyle=\JSONstringstyle,
  morestring=[b]",
  morestring=[d]'
}

\begin{document}

\begin{titlepage}

\newcommand{\HRule}{\rule{\linewidth}{0.5mm}} 

\center 
 
%----------------------------------------------------------------------------------------
%    HEADING SECTIONS
%----------------------------------------------------------------------------------------

\textsc{\LARGE Technical University of Crete}\\[1.5cm]
\IfFileExists{Images/tuc_logo.png}{\includegraphics[scale=0.8]{Images/tuc_logo.png}}{\vspace{0.5cm}}\\[1cm]
\textsc{\Large School of Electrical and Computer Engineering}\\[0.5cm] 
\textsc{\large Big Data Processing and Analysis}\\[0.5cm] 

%----------------------------------------------------------------------------------------
%    TITLE SECTION
%----------------------------------------------------------------------------------------

\HRule \\[0.4cm]
{ \huge \bfseries Sketch-Based vs.\ Counter-Based\\[0.2cm] Heavy Hitters in Streaming Data }\\[0.3cm]
{\large A Comparative Evaluation on Wikimedia Pageviews}\\[0.4cm]
\HRule \\[1.5cm]
 
%----------------------------------------------------------------------------------------
%    AUTHOR SECTION
%----------------------------------------------------------------------------------------

\noindent
\begin{minipage}[t]{0.5\textwidth}
\raggedright\large
\emph{Student:}\\[0.2em]
Alexios Onorios Moumtzis\\
\end{minipage}%
\hfill
\begin{minipage}[t]{0.4\textwidth}
\raggedleft\large
\emph{Instructor:}\\[0.2em]
Minos Garofalakis
\end{minipage}\\[1cm]

%----------------------------------------------------------------------------------------
%    DATE SECTION
%----------------------------------------------------------------------------------------

{\Large July 2026}

\end{titlepage}

\setlength{\headheight}{2cm}
\lhead{\IfFileExists{Images/logo-en-r.png}{\includegraphics[scale=0.6]{Images/logo-en-r.png}}{}}

\tableofcontents
\clearpage


\section{Introduction}

Finding \emph{heavy hitters}, the small set of items that account for a
disproportionate fraction of the traffic in a data stream, is one of the
oldest and most heavily used primitives in streaming data analysis. Since the 
alphabet of potential elements is usually much greater than the memory capacity,
exact methods are not feasible, and there is an extensive range of approximation
methods which require sub-linear space. This range can be categorized into two major categories.:

\begin{itemize}
  \item \textbf{Sketch-based algorithms}, such as the Count--Min Sketch (CMS)
        and the Count Sketch (the hashed variant of the AMS estimator, referred
        to here as \emph{FastAMS}). These maintain a fixed grid of hashed
        counters and estimate the frequency of any queried item.
  \item \textbf{Insert-only / counter-based algorithms}, such as Misra--Gries
        (MG), Space-Saving (SS) and Lossy Counting (LC). These maintain an
        explicit, bounded set of monitored (item, counter) pairs and never store
        more than a fixed number of items.
\end{itemize}

In this project, our aim is to study and compare these two families of algorithms under practical large-scale streaming environments, especially concentrating on
\textbf{accuracy}, \textbf{memory consumption}, \textbf{throughput},
\textbf{scalability across Spark partitions}, and \textbf{distributed
mergeability}. As reference points we additionally include two sampling-based
summaries (Reservoir Sampling and Concise Sampling) and an exact top-$k$
baseline computed with Spark. The goal is deliberately \emph{not} to crown a
single ``best'' algorithm, but to characterise the operational constraints under
which each family becomes preferable.

The workload is drawn from the Wikimedia pageview dataset, a real record of
Wikipedia page requests whose access frequencies follow a strongly skewed,
Zipfian-like distribution. Therefore, it is quite suited to heavy hitter
detection, as heavy hitters are nothing but the pages that get the
maximum hits. In order to test its robustness in controlled settings, we
generate Zipfian data sets with different levels of skew.

All algorithms are implemented in Scala on top of Apache Spark and evaluated
through five benchmark suites. The remainder of this document describes the
algorithms (Section~\ref{sec:algorithms}), the dataset and streaming model
(Section~\ref{sec:dataset}), the experimental design
(Section~\ref{sec:design}), the results (Section~\ref{sec:results}), a
cross-cutting discussion (Section~\ref{sec:discussion}), the threats to validity
(Section~\ref{sec:threats}), and the conclusions (Section~\ref{sec:conclusion}).

\section{Algorithms}
\label{sec:algorithms}

This section summarises the algorithms evaluated, grouped by family. For a
stream of total weight $N=\sum_i f(i)$, a user support $\varphi$ defines the
heavy-hitter threshold $\varphi N$: an item $x$ is a true heavy hitter iff its
true frequency $f(x) \ge \varphi N$. Throughout, $\hat f(x)$ denotes an
algorithm's estimate of $f(x)$.

\subsection{Sketch-based algorithms}

\paragraph{Count--Min Sketch (CMS).}
CMS maintains a $d \times w$ grid of counters with one independent hash function
per row. Each update adds the item's weight to one counter per row, and the
estimate is the \emph{minimum} across rows. CMS never under-counts, so its error
is one-sided: $f(x) \le \hat f(x) \le f(x) + \varepsilon N$ with probability
$1-\delta$, using $w=\lceil e/\varepsilon\rceil$ and $d=\lceil\ln(1/\delta)\rceil$.
Because the error scales with the $L_1$ mass $N$, its space is
$O\!\left(\tfrac{1}{\varepsilon}\log\tfrac{1}{\delta}\right)$, i.e.\ linear in
$1/\varepsilon$.

\paragraph{Count Sketch / FastAMS.}
FastAMS hashes each item to one bucket per row and adds a $\pm1$-signed multiple
of its weight; the estimate is the \emph{median} across rows of the
sign-corrected bucket. It is unbiased but signed (it may over- or under-count),
with error bounded by $\varepsilon\lVert f\rVert_2$. Because the error scales
with the $L_2$ mass, its space is
$O\!\left(\tfrac{1}{\varepsilon^2}\log\tfrac{1}{\delta}\right)$, i.e.\ quadratic
in $1/\varepsilon$, but with a much smaller per-item error on skewed data.

\subsection{Counter-based (insert-only) algorithms}

\paragraph{Misra--Gries (MG).}
MG maintains at most $k-1$ (item, counter) pairs. A hit on a monitored item
increments its counter; otherwise a free slot is used, or all counters are
decremented. MG never over-counts:
$\hat f(x) \le f(x) \le \hat f(x) + \varepsilon N$ with
$k = \lceil 1/\varepsilon\rceil$. Any item with $f(x) > N/k$ is guaranteed to be
retained. Crucially, MG summaries are \emph{mergeable} (with a lossy prune step),
which is central to the distributed evaluation.

\paragraph{Space-Saving (SS).}
SS keeps $m$ monitored counters, each with an over-estimation term
$\varepsilon_i$. On a miss, the least element is evicted and its count reused:
$\text{count}_i - \varepsilon_i \le f(e_i) \le \text{count}_i$. SS therefore
\emph{over}-estimates and is known to be isomorphic to MG (it maintains an upper
rather than a lower bound).

\paragraph{Lossy Counting (LC).}
LC is a deterministic, insert-only summary that divides the stream into buckets
of width $1/\varepsilon$ and prunes at bucket boundaries, guaranteeing
$\hat f(x) \le f(x) \le \hat f(x) + \varepsilon N$. On skewed data the top items
enter in the first bucket and are never pruned, so they tend to be counted
essentially exactly.

\subsection{Sampling baselines and exact reference}

\paragraph{Reservoir Sampling} maintains a fixed-size uniform sample of $M$
stream positions (Vitter's Algorithm R) and estimates
$\hat f(x) = (c_{\text{sample}}(x)/m_{\text{eff}})\cdot N$.
\paragraph{Concise Sampling} keeps a bounded multiset sample with an adaptive
threshold $T$, estimating $\hat f(x) = \text{count}_R(x)\cdot T$.
\paragraph{Exact baseline.} The ground truth is computed with Spark
\texttt{groupBy(key).sum(views)} and stored under \texttt{results/exact\_topk/};
it is used to compute all precision, recall and per-item error figures.

\section{Dataset and Streaming Model}
\label{sec:dataset}

\subsection{Wikimedia pageviews}

The primary workload is the Wikimedia pageview dataset. Raw hourly dumps
(bzip2-compressed) are parsed into a Parquet dataset of \texttt{(key, views)}
pairs, where \texttt{key} concatenates the project/domain and page title and
\texttt{views} is the request count. After cleaning (dropping null counts and
placeholder titles), the resulting stream contains
\textbf{10{,}897{,}474} distinct \texttt{(project, page)} rows with a total
weight of \textbf{$N = 44{,}367{,}267$} views. The frequency distribution is
strongly Zipfian, with $\lVert f\rVert_2 \approx 5.5\times10^6$; at support
$\varphi = 10^{-4}$ the heavy-hitter threshold is $\varphi N = 4{,}437$ views,
yielding \textbf{158} true heavy hitters.

\subsection{Streaming model}

Each Parquet row is replayed as one weighted stream element in the original scan
order via \texttt{df.toLocalIterator()}, i.e.\ a single-pass
\emph{cash-register} model in which every element carries a weight equal to its
view count. This preserves streaming (incremental, one-pass) semantics while
avoiding a full batch re-aggregation inside each algorithm.

There is some importance in clarifying the characteristics of the stream. Since the source dump is already pre-aggregated on the page level for a snapshot window, the workload is best defined as a \emph{weighted single-window stream} rather than a strict
event-by-event replay in timestamp order. This is a faithful representation of the heavy hitter problem under consideration (each row corresponds to an incremental weighted update), but it doesn’t capture any notion of temporal drift between multiple time frames. This is revisited in Section~\ref{sec:threats}.

\subsection{Synthetic Zipfian streams}

To evaluate robustness under controlled skew, synthetic Zipfian streams are
generated with exponent $s \in \{0.6, 0.8, 1.0, 1.2, 1.5\}$ over a fixed
vocabulary, with deterministic keys so that the true top-$k$ is known exactly.
Lower $s$ yields a flatter (more uniform) distribution and higher $s$ a more
concentrated one. Because counts are rounded, the realised total weight varies
with $s$ (from $\approx 1.3\times10^7$ at $s=0.6$ to $\approx 4.4\times10^7$ at
$s=1.5$); the skew evaluation therefore uses a distribution-shape-independent,
threshold-free top-$k$ overlap metric (Section~\ref{sec:design}).

\section{Experimental Design}
\label{sec:design}

\subsection{Spark execution model}

All experiments run on Apache Spark in \texttt{local[*]} mode, using every
available core as a parallel worker thread. The single-structure benchmarks
stream the cached, pre-aggregated \texttt{(key, views)} dataset on the driver.
The distributed benchmark instead builds one summary per Spark partition in
parallel with \texttt{mapPartitions} and combines them with the structure's own
\texttt{merge} via \texttt{treeReduce}; the partition count $P$ is varied over
$\{1,2,4,8,16\}$ using \texttt{df.repartition(P)}. This exercises the same
logical distributed dataflow: partitioned build followed by a tree of merges
that a multi-node cluster would use. It does not model HDFS I/O, network
shuffle cost or executor-level isolation; the implications are discussed in
Section~\ref{sec:threats}.

\subsection{Metrics}

For threshold-based detection, the reported heavy-hitter set is compared against
the true set to obtain true positives (TP), false positives (FP) and false
negatives (FN), from which
$\text{precision} = \text{TP}/(\text{TP}+\text{FP})$ and
$\text{recall} = \text{TP}/(\text{TP}+\text{FN})$ are computed. The relative
frequency error is averaged over the correctly detected heavy hitters,
$\lvert \hat f(x) - f(x)\rvert / f(x)$. Throughput is the number of stream rows
processed per second, and memory is the estimated size of the summary structure.
The skew evaluation is threshold-free: each algorithm produces its own top-$k$
($k=100$) and is scored by overlap with the exact top-$k$, from which
precision, recall and $F_1$ coincide.

\subsection{Benchmark suites}

Five suites are used:
\begin{enumerate}
  \item \textbf{Baseline detection}: each algorithm at literature-tuned
        parameters on the real stream at $\varphi = 10^{-4}$.
  \item \textbf{Same-memory sweep}: all algorithms constrained to identical
        memory budgets (128\,KB--8\,MB) to compare accuracy at equal space.
  \item \textbf{Same-threshold sweep}: the same budgets under a fixed,
        strict detection threshold $\varphi N$.
  \item \textbf{Skew robustness}: fixed memory (1\,MB), fixed top-$k$, over
        the synthetic Zipf sweep plus the real snapshot.
  \item \textbf{Distributed partitioning \& mergeability}: fixed memory,
        varying $P$, measuring scalability and merge fidelity.
\end{enumerate}

\section{Results}
\label{sec:results}

\subsection{Baseline detection on real data}

Table~\ref{tab:baseline} reports each algorithm at its own tuned parameters.
These configurations do \emph{not} share a memory budget: CMS in particular
is given a very large grid, so the table shows what each family can achieve
when tuned independently, not an equal-space comparison (that is the purpose of
Section~\ref{sec:samemem}).

\begin{table}[H]
\centering
\small
\begin{tabular}{llrrrrr}
\hline
\textbf{Algorithm} & \textbf{Family} & \textbf{Memory} & \textbf{Thr.\ (rows/s)} & \textbf{Prec.} & \textbf{Rec.} & \textbf{Mean rel.\ err.} \\
\hline
Count--Min          & Sketch   & 148 MB & 196{,}132 & 98.1\% & 98.1\% & 0.02\% \\
FastAMS (Count Sk.) & Sketch   & 11 MB  & 78{,}109  & 98.1\% & 98.1\% & 0.14\% \\
Misra--Gries        & Counter  & 8.6 MB & 267{,}600 & 92.8\% & 98.1\% & 1.11\% \\
Space-Saving        & Counter  & 13.9 MB& 145{,}786 & 96.3\% & 98.1\% & 1.10\% \\
Lossy Counting      & Counter  & 7.6 MB & 328{,}208 & 90.6\% & 98.1\% & 0.00\% \\
Concise Sampling    & Sampling & 5.8 MB & 351{,}929 & 77.1\% & 93.7\% & 16.57\% \\
Reservoir Sampling  & Sampling & 7.4 MB & 442{,}106 & 75.5\% & 89.9\% & 13.24\% \\
\hline
\end{tabular}
\caption{Baseline heavy-hitter detection on the real Wikimedia stream at
$\varphi=10^{-4}$ (158 true heavy hitters), each algorithm at its own tuned
parameters. Memory budgets differ across rows.}
\label{tab:baseline}
\end{table}

Two observations stand out. First, all summaries recover recall around 98\%; the
family differences show up in \emph{precision} and \emph{per-item error}. The
counter-based summaries and sampling summaries report a few extra candidates
(the MG and LC variants use a recall-preserving report threshold
$(\varphi-\varepsilon)N$, which slightly inflates false positives). Second, LC's
per-item error is effectively zero: on this skewed stream the true heavy hitters
enter early and are never pruned, so they are counted exactly. The sampling
baselines are the fastest but the least accurate, with double-digit relative
error, confirming that uniform sampling is a poor frequency estimator for heavy
hitters.

\begin{figure}[H]
\centering
\includegraphics[width=\linewidth]{Images/fig_baseline.pdf}
\caption{Baseline detection at tuned parameters ($\varphi=10^{-4}$, 158 true heavy hitters). (a) Precision bars per algorithm with recall shown as a dashed line; recall is nearly identical across algorithms except for the two sampling methods. (b) Processing throughput: sampling summaries are fastest, FastAMS (Count Sketch) is slowest at its tuned 11\,MB grid.}
\label{fig:baseline}
\end{figure}

\subsection{Accuracy at equal memory}
\label{sec:samemem}

Constraining every algorithm to the same memory budget exposes the fundamental
space--accuracy trade-off. Table~\ref{tab:samethresh} fixes the budget at
1\,MB under a strict $\varphi N$ threshold.

\begin{table}[H]
\centering
\small
\begin{tabular}{lrrrrr}
\hline
\textbf{Algorithm} & \textbf{Prec.} & \textbf{Rec.} & \textbf{Mean rel.\ err.} & \textbf{Max rel.\ err.} & \textbf{Thr.\ (rows/s)} \\
\hline
Count--Min         & 74.2\% & 100.0\% & 7.47\%  & 19.70\% & 469{,}274 \\
FastAMS            & 99.4\% & 100.0\% & 0.70\%  & 4.73\%  & 129{,}542 \\
Misra--Gries       & 100.0\%& 79.7\%  & 11.21\% & 32.77\% & 280{,}913 \\
Space-Saving       & 58.1\% & 100.0\% & 16.68\% & 66.89\% & 167{,}057 \\
Lossy Counting     & 100.0\%& 98.7\%  & 0.00\%  & 0.00\%  & 315{,}430 \\
Concise Sampling   & 53.8\% & 88.6\%  & 25.90\% & 151.72\%& 411{,}956 \\
Reservoir Sampling & 79.2\% & 84.2\%  & 21.26\% & 110.43\%& 327{,}271 \\
\hline
\end{tabular}
\caption{Equal-memory comparison at a 1\,MB budget under the strict $\varphi N$
detection rule (same-threshold suite).}
\label{tab:samethresh}
\end{table}

At equal memory the picture reverses relative to the tuned baseline: FastAMS and
Lossy Counting are the strongest, combining near-perfect precision and recall
with very small error, while CMS, so dominant when given 148\,MB, now
trades precision for space because its error scales with $N$. The one-sided
error structure of each family is visible directly in the two guaranteed
algorithms: CMS keeps recall at 100\% but sheds precision as space shrinks,
whereas MG keeps precision at 100\% but sheds recall. Table~\ref{tab:sweep}
makes this dual behaviour explicit across the full memory range.

\begin{table}[H]
\centering
\small
\begin{tabular}{lrrr|rrr}
\hline
& \multicolumn{3}{c|}{\textbf{Count--Min}} & \multicolumn{3}{c}{\textbf{Misra--Gries}} \\
\textbf{Mem (KB)} & \textbf{Prec.} & \textbf{Rec.} & \textbf{Rel.\ err.} & \textbf{Prec.} & \textbf{Rec.} & \textbf{Rel.\ err.} \\
\hline
128  & 0.0\%  & 100\% & 68.14\% & 100\% & 34.2\% & 24.08\% \\
256  & 8.9\%  & 100\% & 33.38\% & 100\% & 48.7\% & 21.29\% \\
512  & 60.3\% & 100\% & 15.67\% & 100\% & 68.4\% & 18.15\% \\
1024 & 74.2\% & 100\% & 7.47\%  & 100\% & 79.7\% & 11.21\% \\
2048 & 93.5\% & 100\% & 3.40\%  & 100\% & 90.5\% & 6.47\%  \\
4096 & 97.5\% & 100\% & 1.56\%  & 100\% & 96.2\% & 3.47\%  \\
8192 & 99.4\% & 100\% & 0.71\%  & 100\% & 98.1\% & 1.71\%  \\
\hline
\end{tabular}
\caption{Memory sweep under the strict threshold. CMS never under-counts (recall
fixed at 100\%, precision grows with memory); MG never over-counts (precision
fixed at 100\%, recall grows with memory).}
\label{tab:sweep}
\end{table}

\begin{figure}[H]
\centering
\includegraphics[width=0.75\linewidth]{Images/fig_f1_memory.pdf}
\caption{$F_1$ score across memory budgets (strict $\varphi N$ threshold, same-threshold suite). \textbf{Solid} lines = sketch-based; \textbf{dashed} = counter-based; \textbf{dotted} = sampling. FastAMS and LC reach near-perfect $F_1$ with 1\,MB; CMS needs larger grids to control false positives; SS has an inflection around 1\,MB where its capacity first covers all heavy hitters; sampling baselines lag throughout.}
\label{fig:f1_memory}
\end{figure}

Figure~\ref{fig:f1_memory} extends this picture to the full memory range and to all algorithms.
\textbf{FastAMS} reaches near-perfect $F_1$ already at 512\,KB and saturates at 100\% from 2\,MB
upwards: its $L_2$-bounded error is negligible on this skewed stream even at small budgets.
\textbf{Lossy Counting} climbs steeply between 256\,KB and 1\,MB and also saturates at 100\%
from 2\,MB: at small budgets the capacity is too tight to retain all 158 heavy hitters before
the first prune, but once it is sufficient they are counted exactly (zero relative error).
\textbf{Space-Saving} shows an unusual dip in $F_1$ between 256\,KB and 512\,KB: within this range
the buffer space is saturated with invalid elements that each satisfy the stringent $\varphi N$ criterion,
thus masking real heavy hitters; the inflection point at 1\,MB occurs when the budget becomes
sufficiently large to store all real heavy hitters together.
\textbf{Concise Sampling} and \textbf{Reservoir Sampling} consistently increase but fall behind the
deterministic algorithms for any budget size: the $F_1$ score of both sampling algorithms is limited
by sampling variance and never achieves deterministic performance of the corresponding sketch and 
counter families within the tested range.

\subsection{Robustness to skew}

Table~\ref{tab:skew} reports the top-100 $F_1$ (equal to precision and recall
under fixed $k$) at a 1\,MB budget across the synthetic Zipf sweep and the real
snapshot.

\begin{table}[H]
\centering
\small
\begin{tabular}{lrrrrrr}
\hline
\textbf{Algorithm} & \textbf{$s{=}0.6$} & \textbf{$0.8$} & \textbf{$1.0$} & \textbf{$1.2$} & \textbf{$1.5$} & \textbf{wiki} \\
\hline
Count--Min         & 98 & 99  & 100 & 100 & 100 & 99  \\
FastAMS            & 98 & 100 & 100 & 100 & 100 & 100 \\
Misra--Gries       & 93 & 97  & 99  & 100 & 100 & 96  \\
Space-Saving       & 92 & 97  & 99  & 100 & 100 & 95  \\
Lossy Counting     & 100& 100 & 100 & 100 & 100 & 100 \\
Concise Sampling   & 81 & 88  & 92  & 94  & 91  & 81  \\
Reservoir Sampling & 81 & 93  & 96  & 96  & 95  & 86  \\
\hline
\end{tabular}
\caption{Top-100 $F_1$ (\%) versus Zipf exponent $s$ at a 1\,MB budget. Higher
skew makes the heavy hitters easier to separate.}
\label{tab:skew}
\end{table}

Every algorithm improves monotonically as skew increases, since a more
concentrated head is easier to separate from the tail. The counter-based
summaries and the deterministic LC are the most robust across the whole range;
the sampling baselines degrade sharply on the flatter ($s=0.6$) distributions
and on the real data, where the effective skew is milder than the synthetic
$s=1.5$ case. The accuracy gap between families is small for very skewed distributions but grows considerably as the distribution becomes less skewed.

\begin{figure}[H]
\centering
\includegraphics[width=\linewidth]{Images/fig_skew.pdf}
\caption{Top-100 $F_1$ vs.\ Zipf exponent at 1\,MB budget. Synthetic streams (left of dashed separator) and real Wikimedia snapshot (right). All algorithms improve as skew increases; the gap between families is largest at $s{=}0.6$ and narrows as the distribution concentrates. LC and CMS/FastAMS are robust across all conditions; sampling baselines degrade sharply at low skew.}
\label{fig:skew}
\end{figure}

\subsection{Scalability and distributed mergeability}

Table~\ref{tab:dist} reports the distributed suite at a 1\,MB budget as the
partition count $P$ grows. \texttt{max abs.\ diff.} is the largest deviation of
any top-100 estimate from the single-partition ($P{=}1$) result, and directly
measures merge fidelity.

\begin{table}[H]
\centering
\small
\begin{tabular}{lrrrrr}
\hline
\textbf{Algorithm} & \textbf{$P$} & \textbf{Build+merge (s)} & \textbf{Speedup} & \textbf{Mean rel.\ err.} & \textbf{Max abs.\ diff.\ vs $P{=}1$} \\
\hline
Count--Min   & 1  & 23.18 & 1.00 & 4.35\% & 0 \\
             & 2  & 13.12 & 1.77 & 4.35\% & 0 \\
             & 4  & 7.00  & 3.31 & 4.35\% & 0 \\
             & 8  & 5.22  & 4.44 & 4.35\% & 0 \\
             & 16 & 5.28  & 4.39 & 4.35\% & 0 \\
\hline
FastAMS      & 1  & 51.79 & 1.00 & 0.39\% & 0 \\
             & 2  & 25.89 & 2.00 & 0.39\% & 0 \\
             & 4  & 17.13 & 3.02 & 0.39\% & 0 \\
             & 8  & 13.04 & 3.97 & 0.39\% & 0 \\
             & 16 & 14.04 & 3.69 & 0.39\% & 0 \\
\hline
Misra--Gries & 1  & 30.41 & 1.00 & 9.98\% & 0 \\
             & 2  & 13.40 & 2.27 & 4.87\% & 1{,}777 \\
             & 4  & 7.21  & 4.22 & 2.52\% & 2{,}425 \\
             & 8  & 5.06  & 6.01 & 1.53\% & 2{,}712 \\
             & 16 & 5.15  & 5.91 & 1.20\% & 2{,}823 \\
\hline
\end{tabular}
\caption{Partition scalability and merge fidelity at a 1\,MB budget. Linear
sketches merge losslessly (zero deviation); the Misra--Gries merge is lossy and
drifts from the single-partition result as $P$ grows.}
\label{tab:dist}
\end{table}

This is the central systems result of the project. The two linear sketches (CMS
and FastAMS) merge by summing counter tables, so the merged summary is
bit-identical to the single-pass result regardless of $P$: the deviation column
is exactly zero and accuracy is constant across all partition counts.
Misra--Gries, by contrast, is mergeable only with a lossy prune at each merge, so
its top-$k$ estimates drift from the $P{=}1$ answer, and the deviation grows with
$P$ (up to 2{,}823 views at $P{=}16$). Interestingly, MG's error \emph{against
the ground truth} improves with $P$ (from 9.98\% to 1.20\%): with more
partitions each partition summarises less data more accurately, and the
transient union before pruning carries more information than a single
capacity-bounded pass. All three algorithms scale well up to eight partitions
and then plateau, which is the expected behaviour of a shared-memory
\texttt{local[*]} executor once the cores are saturated.

\begin{figure}[H]
\centering
\includegraphics[width=\linewidth]{Images/fig_distributed.pdf}
\caption{Distributed benchmark results at 1\,MB. (a) Parallel speedup: all three algorithms scale near-linearly to $P{=}8$ then plateau as the shared-memory executor saturates. (b) Merge fidelity: linear sketches merge by summing counter arrays and accumulate \emph{zero} deviation regardless of~$P$; MG's prune-at-merge is lossy and drifts up to 2{,}823 page views at $P{=}16$.}
\label{fig:distributed}
\end{figure}

\section{Discussion: which family, and when}
\label{sec:discussion}

The experiments point towards a more complex, constraint-based response.

\begin{itemize}
  \item \textbf{Lots of memory, one-way guarantee required.}
When memory usage is not the limiting factor, but there is a requirement
for guaranteed no-false-negative behavior, then \textbf{Count--Min} is appealing: low overhead in updates, maximum throughput among all the sketches, and 100\% recall. 
The downside is that the error is proportional to $N$, which makes it costly to use for precision.
  \item \textbf{Limited memory budget.} With equal space allocation, \textbf{FastAMS} and \textbf{Lossy Counting} stand out:
                FastAMS due to the extremely small $L_2$ error for skewed
                inputs, LC due to exact counting of (few) actual heavy-hitters.
                For a strict no false positive requirement,
                \textbf{Misra-Gries} seems the only possible choice, at the expense
                of recall for tiny memory budgets.
  \item \textbf{Distributed / mergeable setting.} Linear sketches
          are clearly preferred here due to their embarrassingly parallel and lossless merging
          capabilities, which makes them suitable for Spark pipeline with partitioning.
          While counter summaries are also mergeable and good for parallelization,
          their merging is approximate and drifts with the number of partitions.
  \item \textbf{Throughput-driven, Accuracy-tolerant.} Sampling-based summaries are
        the fastest, although they shouldn't be used where precise frequency estimates
        are required; they are fast only when the distribution is highly skewed.
\end{itemize}

In short, the sketch family is preferable when memory is scarce or when lossless
distributed merging matters, while the counter family is preferable when a
one-sided (precision) guarantee and simple, interpretable state are required.

\section{Threats to Validity}
\label{sec:threats}

\paragraph{Streaming model.} The work is a weighted single window stream in
Parquet scan order, not in event replay order by timestamps in multiple windows.
The heavy hitter analysis results will remain the same (every row is an incremental weighted update), but temporal drift and windowed heavy hitters are not applicable.

\paragraph{Local Spark.} The distributed suite runs in \texttt{local[*]} mode.
It faithfully reproduces the two properties under study, merge correctness and
relative parallel scaling, but does not capture HDFS I/O, network shuffle cost
or executor-level GC isolation. Absolute throughput on a real multi-node cluster
would differ; the \emph{relative} algorithmic comparisons and the
lossless-vs-lossy merge distinction are expected to carry over.

\paragraph{Detection-rule differences.} The counter-based summaries use a
recall-preserving report threshold $(\varphi-\varepsilon)N$ in the same-memory
suite and a strict $\varphi N$ threshold in the same-threshold suite; precision
and recall are therefore only directly comparable within a suite, not across
suites.

\section{Conclusion}
\label{sec:conclusion}

This project implemented and evaluated a broad set of heavy-hitter algorithms
spanning the sketch-based and counter-based families, plus sampling and exact
baselines, on a real Wikimedia pageview stream and on controlled synthetic
Zipfian streams. Across five benchmark suites the results confirm the expected
theory and quantify the trade-offs: at equal memory the $L_2$ sketch (FastAMS)
and Lossy Counting are the most accurate, Count--Min trades memory for precision
while guaranteeing recall, Misra--Gries guarantees precision while trading
recall, and all families converge as skew increases. The distributed suite
provides the sharpest distinction: linear sketches merge losslessly and are
ideally suited to partitioned Spark execution, whereas counter summaries merge
only approximately. Rather than a single ``best'' algorithm, the study
characterises the operational regimes: memory budget, required error
guarantee, and distribution, under which each family is preferable.

\end{document}